% =====================================================================
%  SN2 (Supplementary Notes): the annotation and review protocol suite.
%  Fragment \input by supplement.tex under the Supplementary Notes part.
%  No preamble here.
% =====================================================================

\section{The annotation and review protocol suite}
\label{sec:supp-protocol-suite}

The annotation workflow described in the main paper has been consolidated
into a versioned, openly published protocol suite available in the project
repository.\footnote{\url{https://github.com/ForomePlatform/genetic-evidence-model/tree/master/protocols}}
The suite distinguishes what a well-formed annotation or review must contain
from how it is performed, either manually or with AI assistance. It governs
both pilot tracks: the curator-authored manual track and the AI-drafted
track.

A mode-agnostic annotation protocol (\texttt{PROTOCOL.md}) defines the
requirements for a valid annotation, including classes, dimensions, source
anchoring, and flag semantics. Two mode-specific specializations adapt this
protocol for autonomous drafting (\texttt{PROTOCOL\_AUTONOMOUS.md}) and
interactive, curator-in-the-loop work. A separate review protocol
(\texttt{REVIEW\_PROTOCOL.md}, with an interactive variant) governs
assertion-by-assertion review of drafted annotations. A labeling-examples
document records worked cases for recurring judgment calls. These written
protocols are also packaged as executable skills
(\cref{sec:supp-skills}), one for annotation and one for review, so that the
same rules can be applied by a human curator or by an AI assistant under
curator supervision.

\subsection{Manual-track protocol}

In the \emph{manual track}, four of the six publications, Jossin 2017,
Davis 2011, Nelson 1992, and Gupta 2015, were annotated by Vladimir
Seplyarskiy, a researcher in human and population genetics specializing in
human germline mutation and mutational processes. He is an Assistant
Professor at UT~Southwestern and was formerly affiliated with Harvard Medical
School and Brigham and Women's Hospital. This expertise covers the
human-genetics and population-genetics knowledge domains represented in the
corpus.

For each paper, the curator:

\begin{enumerate}
  \item read the PDF in full and marked relevant passages with highlights
        and short callouts;
  \item exported those annotations with a helper script to a structured JSON
        record containing page, offset, highlighted text, and free-text
        callout fields;
  \item constructed a one-row summary table containing high-level judgments
        about knowledge domain, method, credibility, resolution, mode of
        inheritance, and related dimensions; and
  \item mapped the highlights and summary row to the schema, producing a
        YAML file conforming to the \texttt{GeneticEvidence} model.
\end{enumerate}

When this conversion required normalization, the change was recorded inline
as a \texttt{normalization\_note} with
\texttt{type: ai\_normalization}, while preserving the original text. For
example, the free-text value ``localisation and protein expression'' was
mapped to the enumeration values \texttt{LOCALIZATION} and
\texttt{EXISTENCE}. These notes allow later audits to compare the original
and normalized values and support subsequent ontology alignment.

All manual-track annotations were validated against the SHACL schema
(\fpath{schema/genetic_evidence.shacl.ttl}). Violations, such as a missing
always-required dimension or a conditional dimension whose activation
condition held but whose value was absent, had to be resolved before an
annotation was admitted to the corpus.

The four manual annotations are treated as \emph{ground truth} for this
work. They are not a gold standard in the statistical sense, because the
corpus has only one annotator. They are, however, the authoritative reference
against which the AI-drafted track is evaluated.

\subsection{AI-drafted-track protocol}

In the \emph{AI-drafted track}, an AI assistant annotated two publications,
Duerr 2006 and Inouye 2018, directly from their PDFs under the same schema.
The four manual annotations served as worked examples.

For each paper, the assistant received:

\begin{itemize}
  \item the schema (\fpath{schema/dimensions.md} and the SHACL file);
  \item the four manual annotations;
  \item the source PDF; and
  \item the autonomous-annotation skill (\cref{sec:supp-skills}).
\end{itemize}

The assistant produced draft YAML annotations in the same format as the
manual annotations. When it judged a mapping uncertain, or concluded that a
judgment required human review, it recorded the issue inline using one of two
flag types:

\begin{itemize}
  \item \texttt{ai\_uncertainty}, for a mapping judged unreliable and
        referred for a human decision; and
  \item \texttt{ai\_assumption}, for a judgment made without explicit
        textual support and requiring reconsideration at review rather than
        silent acceptance.
\end{itemize}

AI-drafted annotations are not accepted without review. One of the authors,
a curator with expertise in genomic data curation and the evidence model
rather than in human or population genetics, reviewed each annotation
assertion by assertion under the review skill. The review assessed:

\begin{itemize}
  \item faithful representation of the source text;
  \item conformance to the schema, including conditional-activation rules;
        and
  \item the appropriateness of any \texttt{ai\_uncertainty} or
        \texttt{ai\_assumption} flags.
\end{itemize}

Corrections were applied directly to the YAML, and every change and ruling
was logged under \fpath{annotations/reviews/}. \Cref{tab:ai-review} summarizes the
resulting verdicts for each \texttt{GeneticEvidence} item, and
\cref{sec:supp-ai-review} analyzes them qualitatively. In all cases, the \emph{curator}, not
the model, is the authority for the recorded annotations.

\subsection{Source anchoring and validation}

A central epistemic commitment of the workflow is that every
\emph{assertion} must be independently auditable against its source.

\begin{itemize}
  \item Each assertion in a \texttt{GeneticEvidence} item carries a
        \texttt{source\_span} containing a page number and a concise key
        phrase, typically fewer than fifteen words, sufficient to locate the
        passage by search in the PDF.
  \item The phrase-length limit is intended to respect fair-use norms for
        copyrighted publisher material while keeping verification practical.
  \item The SHACL schema treats \texttt{source\_span} as mandatory. An
        annotation with any assertion lacking a source anchor is rejected at
        validation time.
\end{itemize}

Across the six-paper pilot, all 95 assertions carry source spans. This is a
property of the validation requirement, not an observed outcome. When a
verbatim key phrase could not be reconstructed from a public PDF, for example
because a value appeared in a publisher PDF table not included in the
repository, the annotation records this explicitly with a
\texttt{source\_span\_is\_paraphrase} flag and an explanatory note. Future
curators can then refine the anchor against the original file.

\subsection{Provenance and protocol versions}

Two provenance points are important for interpreting the corpus.

First, the pilot corpus described in the main paper was initially annotated
under an earlier procedure, ``protocol~v0.'' The suite described here is its
subsequent, generalized formalization. The pilot annotations remain valid
under the current protocol. Some artifacts, such as slightly different
span-selection conventions in the earliest AI-drafted annotations, are
recorded as forward-looking notes rather than retrofitted edits.
One such artifact is a field-name change. Protocol~v0 used
\texttt{key\_phrase} for the source-span phrase, whereas v1 uses
\texttt{phrase}. The corpus retains both spellings. Before RDF validation,
the YAML-to-RDF transformation
(\fpath{src/python/main/forome/gem/extraction/yaml_to_rdf.py}) normalizes
both to the RDF property \texttt{gem:phrase}.

Second, the AI-drafted annotations were produced with Claude Opus~4.7 for
Duerr and Claude Opus~4.8 for Inouye, and both were reviewed with Claude
Opus~4.8 under the executable skills. These implementation details are
reported for reproducibility and are not central to the scientific
contribution. In all cases, every AI-produced annotation and review verdict
was adjudicated by the human curator, who remains the authority for the
recorded annotations.
